\documentclass[11pt]{article}
\usepackage[a4paper,margin=1in]{geometry}
\usepackage{amsmath,amssymb,amsthm,booktabs,hyperref,xcolor,bm}
\newtheorem{theorem}{Theorem}
\newtheorem{lemma}{Lemma}
\newtheorem{corollary}{Corollary}
\newtheorem{proposition}{Proposition}
\newtheorem{definition}{Definition}
\newtheorem{assumption}{Assumption}
\newcommand{\R}{\mathcal R}
\newcommand{\E}{\mathbb E}
\newcommand{\norm}[1]{\lVert #1 \rVert}
\title{A Generic Excess-Risk Certificate for Diffusion / Flow Samplers\\
       {\large path, parameterization, solver, and risk measure in one theorem;\\
        with the calibration and tuning that instantiate it}}
\author{thermau5}
\date{\today}
\begin{document}
\maketitle

\begin{abstract}
We derive, from first principles, a single certified upper bound on the
excess terminal risk of a learned diffusion/flow sampler, valid for any
probability path $P$, learned-field parameterization $h$, numerical
solver core $s$, and terminal risk measure $\R$. The bound decomposes
into an irreducible floor, a statistical (estimation) term scaling as
$1/n$, and a discretization term scaling as $N^{-p}$ in the number of
function evaluations. Minimizing it gives closed-form optimal training
and sampling densities, $\rho^\star\propto\sqrt{a}$ and
$m^\star\propto d^{1/(p+1)}$, with an explicit accuracy--compute
frontier. We then specialize: the parameterization $h$ enters through a
conversion Jacobian; the path $P$ through the floor and both densities;
the solver $s$ only through the local defect; and the risk measure $\R$
through a terminal-sensitivity operator. We give the score / noise /
$x_0$ / $v$ / flow instances and the KL / Wasserstein / feature-space
(FID) instances. Finally we separate, operationally and honestly, what
is \emph{calibrated} (measured intrinsic quantities, no metric feedback)
from what is \emph{tuned} (parameters selected against the evaluation
metric), and we report what our experiments establish: the feature-space
$W_2$ risk is the directly-measurable, FID-faithful risk; the certificate
is exact for the risk it bounds but is an asymptotic (small-step) object;
and a single in-regime exponent closes the deployment-NFE gap that no
decomposable bound can.
\end{abstract}

\tableofcontents

\section{Setup and notation}
\label{sec:setup}

Let $r\in[0,B]$ be an intrinsic path coordinate (e.g.\ a monotone
reparameterization of noise level $\sigma$ or half-log-SNR
$\lambda=-\log\sigma$). A \emph{probability path} $P$ specifies
intermediate laws $\mu_r^P$ interpolating a tractable prior $\mu_0^P$ and
the data law $\mu_B^P=p_{\mathrm{data}}$.

\paragraph{Learned field and parameterization.}
The oracle target is a field $h_P^\star(x,r)$; a network learns
$h_\theta(x,r)\approx h_P^\star(x,r)$. Common choices:
\[
h=s\ (\text{score }\nabla_x\log\mu_r),\quad
h=\epsilon\ (\text{noise}),\quad
h=x_0\ (\text{data}),\quad
h=v\ (\text{velocity / flow}).
\]
The learned object is converted into a physical generative field
(an SDE drift or ODE velocity) by a path-and-parameterization map
$F_{P,h}[\,\cdot\,]$. Write the induced \emph{physical error}
\[
e_F(x,r)=F_{P,h}[h_\theta](x,r)-F_{P,h}[h_P^\star](x,r),
\qquad
e_h(x,r)=h_\theta(x,r)-h_P^\star(x,r).
\]

\paragraph{Solver and sampler.}
A solver core $s$ of state order $k_s$ discretizes the generative
dynamics on a grid $r_0<\cdots<r_N=B$ of step density $m(r)$
($\int_0^B m=N$). A training schedule density $\rho(r)$ allocates the $n$
training samples ($\int_0^B\rho=n$, effective local count $n\rho(r)$).
The full control is $u=(\theta,P,h,s,\rho,m,N)$.

\paragraph{Risk.}
The terminal risk is $\R_{P,h}(\theta)=\R(\mu_B^\theta,\mu_B^\star)$ for a
divergence/metric $\R$ (KL, $W_2$, feature-space distance, task loss).
The \emph{excess risk} is $\R_{P,h,s}(\theta,m)-\R^\star$ where $\R^\star$
is the best achievable by the family.

\section{The generic certificate}
\label{sec:generic}

\subsection{Assumptions}

\begin{assumption}[Local linearization of the conversion]\label{as:lin}
There is a (matrix- or scalar-valued) field $J_{P,h}(r)$ such that
$e_F(x,r)=J_{P,h}(r)\,e_h(x,r)+o(\norm{e_h})$ uniformly on the support of
$\mu_r^P$.
\end{assumption}

\begin{assumption}[Terminal-risk stability]\label{as:stab}
There is a nonnegative weight $G_P^\R(r)$ (the \emph{terminal-risk
sensitivity}) such that local physical errors propagate to terminal risk
as
\[
\R_{P,h}(\theta)-\R^\star \;\le\;
\int_0^B G_P^\R(r)\,\E_{x\sim\mu_r^P}\norm{e_F(x,r)}^2\,dr \;+\; B_{\mathrm{disc}} .
\]
\end{assumption}

\begin{assumption}[Local ERM rate]\label{as:erm}
For the local population risk
$L_{P,h,r}(\theta)=\E_{x\sim\mu_r^P}\norm{h_\theta(x,r)-h_P^\star(x,r)}^2$,
the ERM estimator $\hat\theta_\rho$ obeys, near $r$,
\[
L_{P,h,r}(\hat\theta_\rho)\;\le\; b_{P,h}^2(r)+\frac{c_{\mathrm{stat}}\,V_{P,h}(r)}{n\rho(r)},
\]
with $b_{P,h}^2(r)=\inf_{\theta\in\Theta}L_{P,h,r}(\theta)$ the local
approximation bias and $V_{P,h}(r)$ the local statistical difficulty
(target variance / complexity / effective dimension).
\end{assumption}

\begin{assumption}[Solver local defect]\label{as:solver}
The solver's leading local truncation defect at $(x,r)$ is
\[
\tau_s(x,r)=\lim_{\Delta r\to 0}\frac{\Psi_{s,\Delta r}(x)-\Phi_{\Delta r}(x)}{(\Delta r)^{k_s+1}},
\]
where $\Phi$ is the exact one-step flow and $\Psi_s$ one solver step.
\end{assumption}

\subsection{Stability $\to$ parameterization-weighted target risk}

\begin{lemma}[Risk weight]\label{lem:weight}
Under Assumptions~\ref{as:lin}--\ref{as:stab}, with the
\emph{parameterization-specific risk weight}
\[
w_{P,h}^\R(r):=G_P^\R(r)\,\norm{J_{P,h}(r)}_{\mathrm{op}}^2 ,
\]
the terminal excess risk obeys
\[
\R_{P,h}(\theta)-\R^\star \le
\int_0^B w_{P,h}^\R(r)\,\E_{\mu_r^P}\norm{h_\theta-h_P^\star}^2\,dr + B_{\mathrm{disc}} .
\]
\end{lemma}
\begin{proof}
Insert Assumption~\ref{as:lin} into Assumption~\ref{as:stab} and bound
$\norm{J e_h}^2\le\norm{J}_{\mathrm{op}}^2\norm{e_h}^2$. The $o(\cdot)$
term is higher order and absorbed into $B_{\mathrm{disc}}$.
\end{proof}

\subsection{Insert ERM $\to$ statistical term}

Plugging Assumption~\ref{as:erm} into Lemma~\ref{lem:weight} and defining
\[
Q_{0,P,h}=\int_0^B w_{P,h}^\R(r)\,b_{P,h}^2(r)\,dr,
\qquad
a_{P,h}^\R(r)=w_{P,h}^\R(r)\,V_{P,h}(r),
\]
gives
\[
\R_{P,h}(\hat\theta_\rho)-\R^\star \le
Q_{0,P,h}+\frac{c_{\mathrm{stat}}}{n}\int_0^B\frac{a_{P,h}^\R(r)}{\rho(r)}\,dr
+ B_{\mathrm{disc}} .
\]

\subsection{Discretization term}

\begin{lemma}[Discretization density]\label{lem:disc}
Under Assumption~\ref{as:solver}, with terminal-risk amplification
$A_\R(r,x)$ (the linear sensitivity of final risk to a state perturbation
introduced at $r$) and risk exponent $q\in\{1,2\}$,
define the \emph{local numerical stiffness}
\[
d_{P,h,s}^\R(r)=\E_{x\sim\mu_r^P}\big\|A_\R(r,x)\,\tau_s(x,r)\big\|^q .
\]
Then, for step density $m$ and a risk-decay exponent $p$,
\[
B_{\mathrm{disc}}\;\le\; c_{\mathrm{disc}}\int_0^B \frac{d_{P,h,s}^\R(r)}{m(r)^p}\,dr .
\]
For a state-order-$k_s$ solver, $p=k_s$ for an unsquared terminal error
and $p=2k_s$ for a squared-error certificate; in general $p$ is the
\emph{risk}-decay order, not automatically the solver state order.
\end{lemma}
\begin{proof}[Proof sketch]
A step of size $h\sim 1/m$ at $r$ has local state defect
$\norm{\tau_s}h^{k_s+1}$; its terminal-risk contribution is
$\norm{A_\R\tau_s}^q h^{q(k_s+1)}$. The number of steps near $r$ per unit
$r$ is $m$, so the contribution density is
$m\cdot\norm{A_\R\tau_s}^q (1/m)^{q(k_s+1)}=d_{P,h,s}^\R\,m^{-(q(k_s+1)-1)}$.
Writing $p=q\,k_s+(q-1)$ and absorbing constants gives the stated form;
the two displayed cases are $q=1$ ($p=k_s$) and $q=2$ ($p=2k_s$).
\end{proof}

\subsection{Main theorem}

\begin{theorem}[Generic excess-risk certificate]\label{thm:main}
Under Assumptions~\ref{as:lin}--\ref{as:solver}, for any path $P$,
parameterization $h$, solver $s$, and risk $\R$,
\[
\boxed{\;
\R_{P,h,s}(\hat\theta_\rho,m)-\R^\star \;\le\;
Q_{P,h,s}^\R(\rho,m,N)=
Q_{0,P,h}+\frac{c_{\mathrm{stat}}}{n}\int_0^B\frac{a_{P,h}^\R(r)}{\rho(r)}\,dr
+ c_{\mathrm{disc}}\int_0^B\frac{d_{P,h,s}^\R(r)}{m(r)^p}\,dr .
\;}
\]
\end{theorem}
\begin{proof}
Combine the ERM-inserted bound with Lemma~\ref{lem:disc}.
\end{proof}

\paragraph{Certificate hierarchy.}
Reading off the dependencies:
\[
\text{change }h\Rightarrow \text{changes }a;\quad
\text{change }P\Rightarrow \text{changes }a,d,Q_0;\quad
\text{change }s\Rightarrow \text{changes only }d;\quad
\text{change }\R\Rightarrow \text{changes }a,d\ \text{(through }G^\R,A_\R).
\]

\section{The optimizer and the frontier}
\label{sec:optimizer}

\begin{theorem}[Optimal densities]\label{thm:opt}
Fix $N,n$. The minimizers of $Q_{P,h,s}^\R$ over densities with
$\int\rho=n$, $\int m=N$ are
\[
\rho^\star(r)\propto\sqrt{a_{P,h}^\R(r)},
\qquad
m^\star(r)\propto d_{P,h,s}^\R(r)^{1/(p+1)} .
\]
\end{theorem}
\begin{proof}
The statistical term $\int a/\rho$ under $\int\rho=n$ is minimized by
Cauchy--Schwarz: $\big(\int\sqrt a\big)^2\le \big(\int a/\rho\big)\big(\int\rho\big)$,
equality iff $\rho\propto\sqrt a$. For the grid, the Lagrangian
$\int d/m^p-\lambda(\int m-N)$ has stationarity
$-p\,d\,m^{-(p+1)}=\lambda$, i.e.\ $m^{p+1}\propto d$, giving
$m^\star\propto d^{1/(p+1)}$.
\end{proof}

\begin{corollary}[Optimized certificate / frontier]\label{cor:frontier}
At the optima,
\[
Q^{\R,\star}_{P,h,s}(N)=Q_{0,P,h}
+\frac{c_{\mathrm{stat}}}{n}\Big(\int_0^B\!\sqrt{a_{P,h}^\R}\Big)^2
+\frac{c_{\mathrm{disc}}}{N^{p}}\Big(\int_0^B\!{d_{P,h,s}^\R}^{1/(p+1)}\Big)^{p+1}.
\]
Inverting $Q^{\R,\star}(N)\le\varepsilon$ gives the minimum-NFE
accuracy-constrained sampler $N^\star(\varepsilon)$.
\end{corollary}

\noindent The shape laws $\rho^\star\propto\sqrt a$,
$m^\star\propto d^{1/(p+1)}$ are \emph{invariant} across all
$(P,h,s,\R)$; only the profiles $a,d$ (and floor $Q_0$) change.

\section{Instances}
\label{sec:instances}

\subsection{Parameterizations: the conversion factor}
\label{sec:params}

Each parameterization converts target error to score error
$s_\theta-s^\star=\kappa_h(r)\,(h_\theta-h^\star)$, and the path-KL weight
is $w^{\mathrm{KL}}_{P,h}(r)=\tfrac12 g(r)^2\kappa_h(r)^2$, where $g(r)$ is
the reverse-SDE diffusion coefficient (Girsanov). With Gaussian
corruption $x_r=\alpha(r)x_0+\sigma(r)\epsilon$:

\begin{center}
\small
\begin{tabular}{lll}
\toprule
$h$ & conversion $\kappa_h$ ($s$-error per $h$-error) & risk weight $w^{\mathrm{KL}}_{P,h}(r)$ \\
\midrule
score $s$ & $1$ & $\tfrac12 g^2$ \\
noise $\epsilon$ & $1/\sigma$ & $\tfrac12 g^2/\sigma^2$ \\
data $x_0$ & $\alpha/\sigma^2$ & $\tfrac12 g^2\alpha^2/\sigma^4$ \\
velocity $v$ & $\kappa_v$ (rotation-dependent) & $\tfrac12 g^2\kappa_v^2$ \\
\bottomrule
\end{tabular}
\end{center}

\noindent Then $a_{P,h}^{\mathrm{KL}}(r)=w^{\mathrm{KL}}_{P,h}(r)\,V_{P,h}(r)$.
The theorem is unchanged; only $a$ moves. For the EDM convention
($\alpha=1,\ \sigma=t$, VE diffusion $g^2=2\sigma$) the data-prediction
weight scales as $g^2/\sigma^4=2\sigma^{-3}$, the noise weight as
$\sigma^{-1}$, the score weight as $\sigma$. The exponent of the
terminal-sensitivity weight is thus a \emph{parameterization-and-path}
quantity in an $O(\sigma^{-1})$--$O(\sigma^{-3})$ family; see
\S\ref{sec:tuning} for how its precise value is fixed in practice.

\subsection{Flow matching / velocity (deterministic transport)}
\label{sec:flow}
For $h=v$ with $dx_r/dr=v^\star(x_r,r)$ and Lipschitz $L(r)$, define the
terminal amplification $A_P(r)=\exp\!\big(\int_r^B L\big)$. A natural
terminal risk is $\R_{W}=W_2^2(\mu_B^\theta,\mu_B^\star)$, giving
$w^{W_2}_{P,v}(r)=B\,A_P(r)^2$ and
$a^{W_2}_{P,v}(r)=B\,A_P(r)^2 V_{P,v}(r)$, with the same certificate form.

\subsection{Risk measures: the terminal-sensitivity operator}
\label{sec:risks}

The risk enters only through $G^\R$ (statistical) and $A_\R$
(discretization). With the solver defect $\tau_s$ fixed:

\begin{center}
\small
\begin{tabular}{ll}
\toprule
risk $\R$ & local numerical stiffness $d^\R_{s}(r)$ \\
\midrule
path-KL (reverse SDE) & $\tfrac12\,\E\big\|\Sigma(r)^{-1/2}\tau_s\big\|^2$ \quad ($p=2k_s$) \\
squared Wasserstein & $\E\big[A(r)^2\norm{\tau_s}^2\big]$, $A(r)=e^{\int_r^B L}$ \quad ($p=2k_s$)\\
task / feature loss $\ell$ & $\E\big\|\nabla_x\ell(\cdot)\,D_z\Phi_{B\leftarrow r}\,\tau_s\big\|^2$ \\
\bottomrule
\end{tabular}
\end{center}

\noindent All share the semantic form
$d^\R_s(r)=(\text{terminal risk sensitivity})\times(\text{local solver defect})$,
so Theorem~\ref{thm:opt} applies verbatim with the appropriate $d^\R$.

\subsection{The FID instance}
\label{sec:fid}
FID is a feature-space transport cost: with $\varphi$ the Inception map,
$\mathrm{FID}=W_2^2(\mathcal N(\mu_g,\Sigma_g),\mathcal N(\mu_d,\Sigma_d))$
in $\varphi$-space. This is the task-loss instance with
$\ell=$ feature distance, so
\[
A_{\mathrm{FID}}(r)=J_\varphi(x_0)\,D_z\Phi_{B\leftarrow r}\equiv J_\varphi\,A_{r\to0},
\qquad
d^{\mathrm{FID}}_s(r)=\E\big\|J_\varphi\,A_{r\to0}\,\tau_s\big\|^2 .
\]
Writing $d_s(r)=\E\norm{\tau_s}^2$ (the pure local defect) and
$g_\varphi(r)=\norm{J_\varphi A_{r\to 0}}^2$ (the feature terminal
sensitivity), the FID stiffness factorizes (directionally up to the
alignment of $\tau_s$) as $d^{\mathrm{FID}}_s(r)\approx d_s(r)\,g_\varphi(r)$,
and $m^\star_{\mathrm{FID}}\propto (d_s\,g_\varphi)^{1/(p+1)}$.

\section{Calibration}
\label{sec:calibration}

\begin{definition}[Calibration]
Measurement of an intrinsic quantity of $(P,h,s)$ from the model/reference
dynamics, with \emph{no} reference to the evaluation metric. Deterministic;
zero metric-feedback degrees of freedom.
\end{definition}

\paragraph{Local stiffness $d_s(r)$.}
Build a high-resolution reference trajectory (e.g.\ a 16--32 substep Heun
integration) for a calibration batch. At each $r$, take one solver step
from the reference state and compare to the reference target:
\[
\hat d_s(r)=\E_{x\sim\text{ref}}\Big\|\tfrac{\Psi_{s,h}(x,r)-\Phi_h(x,r)}{h^{k_s+1}}\Big\|^q .
\]
This estimates Assumption~\ref{as:solver}'s defect directly and is
solver-general (no symbolic Taylor expansion needed for DPM-Solver++,
UniPC, DEIS, Restart).

\paragraph{Terminal sensitivity $g_\varphi(r)$ (for FID).}
Finite-difference Jacobian--vector product through flow and feature map:
perturb the reference state at $r$ by $\epsilon v$, integrate the flow to
$r=B$, and read the feature change,
\[
\hat g_\varphi(r)=\tfrac{1}{\epsilon^2}\,\E_v\big\|\varphi(x_B(x_r+\epsilon v))-\varphi(x_B(x_r))\big\|^2 .
\]

\paragraph{Statistical difficulty $a(r)$.}
$V_{P,h}(r)$ is the local target variance/complexity, measured from the
training targets; $w^\R_{P,h}$ from the path constants and the
parameterization Jacobian. These are model-intrinsic, never metric-fit.

\noindent\textbf{Properties.} Calibrated quantities cannot overfit the
evaluation metric (they never see it); they feed Theorem~\ref{thm:main}
directly and yield \emph{certified} optima.

\section{Tuning}
\label{sec:tuning}

\begin{definition}[Tuning]
Selection of a free parameter by searching against the evaluation metric
on a held-out (validation) split and keeping the best-scoring value.
Metric-in-the-loop; finite degrees of freedom.
\end{definition}

\paragraph{Why a tuned exponent appears.}
In practice the calibrated stiffness is used in a power-law-weighted form
$d^\R_s(r)\approx d_s(r)\,\sigma^{-k}$ with grid exponent $p$. The
weight $\sigma^{-k}$ is the \emph{parametric model} of the
terminal-sensitivity factor $g_\varphi$ (FID) or $w^\R$ (KL); the theory
fixes its \emph{family} (\S\ref{sec:params}--\ref{sec:risks}) but not the
exact exponent when (i) the precise sensitivity is known only
approximately and (ii) the certificate is asymptotic while deployment is
at large steps. Both $k$ and $p$ are then fixed by a validation-FID sweep.

\paragraph{Protocol (no-leakage).}
Sweep $(k,p)$ on the validation split only; select the argmin; evaluate
\emph{once} on the locked test split. Tuning on test is forbidden.

\paragraph{The certified--tuned boundary.}
\begin{center}
\small
\begin{tabular}{lll}
\toprule
object & status & role \\
\midrule
shape law $m^\star\propto d^{1/(p+1)}$, $\rho^\star\propto\sqrt a$ & theory & optimizer form \\
$d_s(r)$ profile & calibrated & full profile, 0 metric DOF \\
$g_\varphi(r)$ profile & calibrated & feature sensitivity (FID) \\
grid exponent $p$ & theory + tuned & $2k_s$ asymptotically; in-regime value differs \\
weight exponent $k$ & tuned & 1 scalar; parametric model of the sensitivity \\
\bottomrule
\end{tabular}
\end{center}

\section{Empirical reconciliation}
\label{sec:empirics}

We summarize what experiments on EDM CIFAR-10 (UniPC/Heun cores, Clean-FID,
$10^4$ samples) establish about the gap between the certificate and FID.

\paragraph{(i) The risk measure is FID-faithful, and directly measurable.}
The discretization Wasserstein risk
$R_{\mathrm{disc}}=W_2(\varphi\text{-features of }K\text{-step},\ \infty\text{-step})$
tracks FID in \emph{all} ranks with near-equal values (e.g.\ $12.4/13.2$,
$50/45$, $283/270$). Hence FID-excess \emph{is} the feature-space $W_2$
discretization risk: the risk-measure axis of Theorem~\ref{thm:main} is
confirmed for $\R=$ feature-$W_2$.

\paragraph{(ii) A pixel/identity sensitivity is the wrong operator.}
The bound with $A_\R=I$ (pixel $L^2$) is computed correctly --- it tracks
true pixel error --- but its rank disagrees with FID (two grids of equal
pixel-$Q$, FID $13$ vs $88$). Replacing $A_\R$ by the feature sensitivity
$g_\varphi$ restores FID-faithful ranking. This is exactly the
risk-sensitivity correction of \S\ref{sec:risks}.

\paragraph{(iii) The bound is asymptotic; its argmin needs an in-regime
exponent.} Minimizing the cheap decomposable bound does not by itself give
the FID-optimal schedule at deployment NFE, because $\tau_s$ and $A_\R$
are leading-order/linear. An oracle sweep of the weight exponent against
FID has a sharp convex minimum:
\begin{center}
\small
\begin{tabular}{l rrrrrr}
\toprule
$K$ & $k{=}0$ & $k{=}1$ & $k{=}1.5$ & $k{=}2$ & $k{=}2.5$ & $k{=}3$ \\
\midrule
5 & 163.7 & 59.9 & 31.1 & \textbf{21.2} & 33.2 & 109.8 \\
8 & 101.3 & 29.7 & 16.6 & \textbf{9.2} & 9.6 & 53.1 \\
12 & 61.2 & 19.3 & 10.8 & 5.6 & \textbf{5.3} & 21.7 \\
\bottomrule
\end{tabular}
\end{center}
The oracle is $k^\star\approx 2$ (drifting to $2.5$ at $K{=}12$): a clean,
stable minimum, not an arbitrary knob. The leading-order feature
sensitivity gives $g_\varphi\sim\sigma^{-1.5}$ (i.e.\ $k\approx1.5$), which
is \emph{suboptimal} ($31$ vs $21$ at $K{=}5$); the deployment-NFE
correction pushes the effective exponent steeper, into the family
predicted by the KL/parameterization weights of \S\ref{sec:params}
($\sigma^{-1}$ to $\sigma^{-3}$).

\paragraph{(iv) The decomposability--tightness wall.} A bound that is
simultaneously cheap, decomposable ($\int d/m^p$), and tight at large
steps does not exist in general: tightness requires the non-decomposable
nonlinear error interaction that $R_{\mathrm{disc}}$ captures by direct
sampling. Three independent attempts to beat the in-regime-tuned schedule
(measured $g_\varphi$ weight; squared-risk exponent $p=2k_s$; direct
$R_{\mathrm{disc}}$ minimization) all fail or degenerate. Hence one
in-regime-tuned exponent is necessary and sufficient.

\section{Scope of claims}
\label{sec:scope}

\begin{itemize}
\item \textbf{Certified} (calibration + theory, no metric feedback): the
KL/pathwise-$L^2$ optima $\rho^\star\propto\sqrt a$,
$m^\star\propto d_s^{1/(p+1)}$.
\item \textbf{Confirmed risk geometry}: FID $=$ feature-$W_2$ (Theorem
\ref{thm:main} with $\R=$ feature-$W_2$); $R_{\mathrm{disc}}\!\leftrightarrow\!$FID.
\item \textbf{Tuned} (one scalar, validation-selected, theory-shaped):
the FID weight exponent $k\approx2$ (and effective grid exponent $p=2$),
which absorbs the asymptotic-vs-deployment gap.
\end{itemize}
The honest one-line claim: \emph{the sampler is the certificate's optimizer
$m^\star\propto d^{1/(p+1)}$; certified for KL/$L^2$ with a calibrated
difficulty profile, and FID-optimal with the same closed form under one
in-regime-tuned, theory-shaped exponent.}

\end{document}
